\chapter{Trigonometric and Inverse Functions}

The functions in this chapter were not invented by people doing calculus. They were invented thousands of years before, by astronomers tracking the rotation of the heavens and surveyors dividing up land. Sine and cosine are the mathematics of \emph{going around}: of anything that repeats, oscillates, rotates, or vibrates. A swinging pendulum, an alternating current, a sound wave, the seasons, the orbit of a planet, the position of a point on a spinning wheel -- all of them, when you measure them carefully, turn out to be made of sines and cosines.

So before we can do calculus \emph{with} these functions -- before we can differentiate $\sin x$ or integrate $\cos x$ -- we need to really understand what they \emph{are}. Not as buttons on a calculator, and not as ratios in a right triangle (though we will see that picture too), but as something deeper and more useful: the coordinates of a point traveling around a circle. That single picture will explain every identity, every graph, every symmetry, and -- crucially -- it will tell us why we measure angles in those strange units called \emph{radians}, and why doing so is the secret that makes calculus with these functions come out clean.

We will also meet the \emph{inverse} trig functions, the functions that answer the question ``what angle has this sine?'' Inverses are subtle -- a periodic function repeats, so it cannot be undone without care -- and sorting out exactly how to undo a sine teaches a lesson about inverse functions in general that will pay off later.

Let us build it all from the circle up.

\section{Angles: What Are We Really Measuring?}
\label{ch04:angles}

Start with the most basic question. An angle measures \emph{how much you have turned}. Stand facing north, rotate to face east: you have turned through one quarter of a full turn. Rotate to face south: half a turn. The natural, dimensionless way to talk about turning is as a \emph{fraction of a full revolution}.

Everything after that is just choosing a \emph{unit} -- a name for ``one full turn'' -- and the choice is pure convention. The Babylonians, working in base 60 and tracking a sky that seemed to cycle once a year through roughly 360 days, declared one full turn to be $360$ \emph{degrees}. There is nothing sacred about $360$; it is a culturally inherited number that happens to divide nicely ($360 = 2^3\cdot 3^2 \cdot 5$, so halves, thirds, quarters, fifths, sixths, eighths, ninths, tenths, twelfths all come out whole). It is a fine unit for carpentry and navigation. It is a \emph{terrible} unit for calculus, and by the end of this chapter you will feel in your bones exactly why.

\begin{intuitionbox}
An angle is not a length and not an area. It is a \emph{turn} -- a fraction of a full rotation. Degrees and radians are just two different rulers laid against that same turn. Degrees call a full turn $360$; radians call it $2\pi$. Same turn, different number. The whole game is choosing the ruler that makes the rest of mathematics simplest, and that ruler is radians.
\end{intuitionbox}

\subsection{Radians: Measuring an Angle by the Arc It Cuts}

Here is a more natural idea, one that uses no arbitrary number at all. Draw a circle of radius $r$ and an angle at its center. That angle opens up onto an \emph{arc} -- a piece of the circle's edge. The bigger the angle, the longer the arc. In fact the arc length is exactly proportional to the angle. So why not let the arc \emph{be} the measure of the angle?

There is one wrinkle: a bigger circle gives a longer arc for the same angle, so raw arc length depends on $r$. We fix that by dividing out the radius. Define:
\[
  \theta \;=\; \frac{\text{arc length}}{\text{radius}} \;=\; \frac{s}{r}.
\]
This ratio is the same for every circle, because doubling $r$ doubles $s$. The number you get is the angle measured in \textbf{radians}.

\begin{definition}[Radian]
\label{ch04:def-radian}
The radian measure of a central angle is the length of the arc it subtends divided by the radius of the circle: $\theta = s/r$. Equivalently, an angle of one radian subtends an arc whose length equals the radius.
\end{definition}

Because $s = r\theta$, an angle of one radian is the angle for which the arc is \emph{exactly as long as the radius}. Wrap one radius-length of string along the rim of the circle, and the angle it spans at the center is one radian -- a bit under $60^\circ$ ($57.2958\ldots^\circ$, to be precise).

\begin{figure}[H]
\centering
\begin{tikzpicture}[scale=2.4]
  \draw[->] (-1.3,0) -- (1.4,0) node[right]{$x$};
  \draw[->] (0,-1.3) -- (0,1.4) node[above]{$y$};
  \draw[curveblue,thick] (0,0) circle (1);
  % radius along x-axis
  \draw[thick] (0,0) -- (1,0) node[midway,below]{$r$};
  % angle of 1 radian
  \draw[thick] (0,0) -- ({cos(57.2958)},{sin(57.2958)});
  % the arc, drawn thick, equal to radius
  \draw[tangentred,very thick] (1,0) arc (0:57.2958:1);
  \node[tangentred] at (1.18,0.42) {arc $=r$};
  \draw (0.32,0.06) node[above right] {\small $1$ rad};
  \node[curveblue] at (-0.72,0.78) {radius $r$};
\end{tikzpicture}
\caption{One radian: the central angle whose arc length equals the radius. Since the full circumference is $2\pi r$, a full turn is $2\pi$ radians.}
\label{ch04:fig-radian}
\end{figure}

How many radians in a full turn? The whole circumference has length $2\pi r$, so a full revolution measures $s/r = 2\pi r / r = 2\pi$ radians. That gives us the master conversion:
\[
  2\pi \text{ radians} \;=\; 360^\circ,
  \qquad\text{so}\qquad
  \pi \text{ radians} = 180^\circ,
  \qquad
  1 \text{ rad} = \frac{180^\circ}{\pi}.
\]
A handful of values worth knowing cold:
\[
  30^\circ = \tfrac{\pi}{6},\quad
  45^\circ = \tfrac{\pi}{4},\quad
  60^\circ = \tfrac{\pi}{3},\quad
  90^\circ = \tfrac{\pi}{2},\quad
  180^\circ = \pi,\quad
  360^\circ = 2\pi.
\]

\begin{happeningbox}
When you write ``$\theta = \pi/3$'' you are not really writing an abstract angle -- you are writing a \emph{length}. You are saying: on the unit circle (radius $1$), walk a distance of $\pi/3 \approx 1.047$ along the rim. The angle and the arc length are literally the same number when $r=1$. This is why, throughout calculus, the variable inside $\sin$ and $\cos$ feels like a distance: in radians on the unit circle, it \emph{is} one. Degrees hide this by inserting the arbitrary scaling factor $\pi/180$, which is exactly the factor that will later clutter up every derivative if you let it.
\end{happeningbox}

Two immediate payoffs of the $s = r\theta$ definition, both free:

\textbf{Arc length.} The length of an arc of radius $r$ spanning angle $\theta$ (in radians) is simply $s = r\theta$. No factor of $\pi/180$, no fuss.

\textbf{Sector area.} A full disk has area $\pi r^2$ and spans $2\pi$ radians. A sector of angle $\theta$ is the fraction $\theta/(2\pi)$ of the disk, so its area is
\[
  A \;=\; \frac{\theta}{2\pi}\cdot \pi r^2 \;=\; \tfrac12 r^2 \theta.
\]
Hold onto that formula $A = \tfrac12 r^2\theta$; we will use it in a clever sandwich argument later in the book to pin down the most important limit in trigonometry.

\begin{walkbox}
A bicycle wheel has radius $0.35$ m. As you ride, the wheel turns. Convince yourself that the distance the bike travels equals the arc length traced on the rim, and hence that if the wheel turns through $\theta$ radians the bike moves $0.35\,\theta$ meters. Now feel why ``radians'' is the natural rotational unit for anything that rolls: the angle, scaled by the radius, \emph{is} the distance. What ugly conversion factor would appear if you insisted on degrees?
\end{walkbox}

\section{Sine and Cosine as Coordinates on the Unit Circle}
\label{ch04:unit-circle}

Now the central idea of the whole chapter. Forget triangles for a moment. Draw the \textbf{unit circle}: the circle of radius $1$ centered at the origin, the set of points with $x^2 + y^2 = 1$.

Start at the point $(1,0)$ -- three o'clock -- and walk counterclockwise along the rim. After walking an arc-length $t$ (equivalently, after turning through an angle of $t$ radians, since $r=1$), you arrive at some point on the circle. \emph{That point's coordinates are the definition of cosine and sine.}

\begin{definition}[Sine and cosine, circle definition]
\label{ch04:def-sincos}
Let $t$ be a real number. Travel counterclockwise a distance $t$ along the unit circle, starting from $(1,0)$. The coordinates of the point you reach are denoted $(\cos t,\ \sin t)$. That is, $\cos t$ is the horizontal coordinate and $\sin t$ is the vertical coordinate of the point at arc-length $t$ on the unit circle.
\end{definition}

\begin{figure}[H]
\centering
\begin{tikzpicture}[scale=2.6]
  \draw[->] (-1.35,0) -- (1.45,0) node[right]{$x$};
  \draw[->] (0,-1.35) -- (0,1.45) node[above]{$y$};
  \draw[curveblue,thick] (0,0) circle (1);
  \def\ang{52}
  % the arc traversed
  \draw[tangentred,very thick] (1,0) arc (0:\ang:1);
  % radius to the point
  \draw[thick] (0,0) -- (\ang:1);
  \filldraw[black] (\ang:1) circle (0.5pt);
  % drop perpendiculars
  \draw[dashed,targetgreen] (\ang:1) -- ({cos(\ang)},0);
  \draw[dashed,structurepurple] (\ang:1) -- (0,{sin(\ang)});
  % labels
  \node[above right] at (\ang:1) {$(\cos t,\ \sin t)$};
  \node[targetgreen,below] at ({cos(\ang)},0) {$\cos t$};
  \node[structurepurple,left] at (0,{sin(\ang)/2}) {$\sin t$};
  \node[tangentred] at (1.16,0.36) {arc $=t$};
  \draw (0.28,0.05) node[above right] {\small $t$};
\end{tikzpicture}
\caption{Cosine and sine are the coordinates of the point reached by walking arc-length $t$ counterclockwise around the unit circle from $(1,0)$. Cosine is horizontal, sine is vertical.}
\label{ch04:fig-unitcircle}
\end{figure}

This definition beats the old ``opposite over hypotenuse'' from school in three ways. It works for \emph{every} real number $t$, not just angles between $0$ and $90^\circ$: you can keep walking past the top of the circle, past the left side, all the way around and around. It makes the symmetries obvious, as the next section shows. And it ties $\sin$ and $\cos$ directly to motion along a curve, which is exactly what derivatives are about -- the payoff that matters for calculus.

The connection to the right-triangle picture is immediate. Look at Figure~\ref{ch04:fig-unitcircle}: the radius to the point, the dashed vertical segment, and the segment along the $x$-axis form a right triangle with hypotenuse $1$. The horizontal leg is $\cos t$, the vertical leg is $\sin t$. So the familiar SOH-CAH-TOA ratios are just this picture scaled: in a right triangle with hypotenuse $h$, the side opposite angle $t$ is $h\sin t$ and the side adjacent is $h\cos t$. The unit circle is the master triangle, hypotenuse normalized to $1$.

\begin{intuitionbox}
Picture a point riding counterclockwise around the unit circle at steady speed. Its \emph{shadow on the $x$-axis} -- the horizontal coordinate -- traces out $\cos t$ as $t$ increases. Its \emph{shadow on the $y$-axis} traces out $\sin t$. Both shadows slosh back and forth between $-1$ and $1$ forever. Sine and cosine are nothing more than the two shadows of uniform circular motion. Every wave you will ever meet is a shadow of something going around.
\end{intuitionbox}

\subsection{Everything That Follows Immediately}

From this one definition, a cascade of facts falls out with almost no work. This is the sign of a good definition.

\textbf{The Pythagorean identity.} The point $(\cos t, \sin t)$ lies \emph{on} the unit circle, and the unit circle is the set with $x^2 + y^2 = 1$. Therefore, for every real $t$,
\[
  \boxed{\;\cos^2 t + \sin^2 t = 1.\;}
\]
This is not a theorem you prove with effort; it \emph{is} the equation of the circle, with $x$ renamed $\cos t$ and $y$ renamed $\sin t$. (The notation $\cos^2 t$ is shorthand for $(\cos t)^2$, an old convention you simply have to memorize.)

\textbf{Range.} Both coordinates of a point on the unit circle lie between $-1$ and $1$. So $-1 \le \sin t \le 1$ and $-1 \le \cos t \le 1$ for all $t$. The functions never escape that band.

\textbf{Periodicity.} Walking an extra full lap -- an extra $2\pi$ of arc -- brings you back to the same point. Hence
\[
  \sin(t + 2\pi) = \sin t,\qquad \cos(t + 2\pi) = \cos t.
\]
The functions repeat with \textbf{period} $2\pi$. This is why they model anything cyclic.

\textbf{Special values.} Read them straight off the circle. At $t=0$ you are at $(1,0)$, so $\cos 0 = 1,\ \sin 0 = 0$. A quarter turn $t = \pi/2$ lands you at the top $(0,1)$: $\cos(\pi/2) = 0,\ \sin(\pi/2) = 1$. Half a turn $t=\pi$ puts you at $(-1,0)$. Three-quarters $t = 3\pi/2$ at the bottom $(0,-1)$.

\begin{center}
\begin{tabular}{c|ccccc}
\toprule
$t$ & $0$ & $\dfrac{\pi}{6}$ & $\dfrac{\pi}{4}$ & $\dfrac{\pi}{3}$ & $\dfrac{\pi}{2}$ \\[0.6em]
\midrule
$\sin t$ & $0$ & $\dfrac{1}{2}$ & $\dfrac{\sqrt2}{2}$ & $\dfrac{\sqrt3}{2}$ & $1$ \\[0.8em]
$\cos t$ & $1$ & $\dfrac{\sqrt3}{2}$ & $\dfrac{\sqrt2}{2}$ & $\dfrac{1}{2}$ & $0$ \\[0.4em]
\bottomrule
\end{tabular}
\end{center}

Notice the elegant pattern: the sine row reads $\tfrac{\sqrt0}{2}, \tfrac{\sqrt1}{2}, \tfrac{\sqrt2}{2}, \tfrac{\sqrt3}{2}, \tfrac{\sqrt4}{2}$, and the cosine row is the same thing backwards. The values at $\pi/6, \pi/4, \pi/3$ come from the two special triangles -- the $30$-$60$-$90$ and the $45$-$45$-$90$ -- which we will derive as an example below.

\begin{labbox}
Let us nail down $\sin(\pi/4)$ and $\cos(\pi/4)$ from the circle. At $t = \pi/4$ (a $45^\circ$ turn) you are exactly on the diagonal line $y = x$, by symmetry: the point is equidistant from the two axes. So $\cos(\pi/4) = \sin(\pi/4)$; call the common value $c$. Pythagoras says $c^2 + c^2 = 1$, hence $2c^2 = 1$, so $c = 1/\sqrt2 = \sqrt2/2 \approx 0.707$. You just computed a ``table value'' with nothing but the circle and Pythagoras -- no calculator, no memorization.
\end{labbox}

\section{The Other Four: Tangent and Friends}
\label{ch04:tangent}

Sine and cosine are the parents; the remaining four trig functions are built from them by division. The most important is the \textbf{tangent}:
\[
  \tan t \;=\; \frac{\sin t}{\cos t}.
\]
Geometrically, $\tan t$ is the \emph{slope} of the radius line to the point $(\cos t,\sin t)$ -- it is rise over run, $\sin t / \cos t$. It is also, beautifully, the length cut off on the vertical tangent line to the circle at $(1,0)$, which is where the name comes from (\emph{tangent} = touching). The remaining three are reciprocals, named for historical convenience:
\[
  \cot t = \frac{\cos t}{\sin t} = \frac{1}{\tan t},\qquad
  \sec t = \frac{1}{\cos t},\qquad
  \csc t = \frac{1}{\sin t}.
\]
Each is undefined wherever its denominator hits zero. Tangent and secant blow up where $\cos t = 0$, i.e.\ at $t = \pm\pi/2, \pm 3\pi/2,\ldots$; cotangent and cosecant blow up where $\sin t = 0$, i.e.\ at $t = 0, \pm\pi, \pm 2\pi,\ldots$ Dividing the Pythagorean identity $\cos^2 + \sin^2 = 1$ through by $\cos^2 t$ gives a useful cousin:
\[
  1 + \tan^2 t = \sec^2 t,
\]
and dividing instead by $\sin^2 t$ gives $1 + \cot^2 t = \csc^2 t$. These are the same circle equation wearing different clothes.

\begin{trapbox}
$\tan t$ has period $\pi$, \emph{not} $2\pi$. After half a lap the point $(\cos t,\sin t)$ flips to its antipode $(-\cos t, -\sin t)$, and the two sign flips cancel in the ratio $\sin t/\cos t$. So $\tan(t+\pi) = \tan t$. Students who blindly assume every trig function repeats every $2\pi$ get tangent wrong. Always reason from the circle, not from a memorized rule.
\end{trapbox}

\section{Symmetries: Reading Identities Off the Circle}
\label{ch04:identities}

A common student fear is that trigonometry is a swamp of identities to memorize. It is not. Almost every identity is a geometric symmetry of the circle, \emph{seen} rather than recalled. Here are the main ones, each as a picture.

\textbf{Even and odd.} Replace $t$ by $-t$: instead of walking counterclockwise you walk clockwise the same distance. You land at the mirror image across the $x$-axis. Mirroring across the $x$-axis keeps the horizontal coordinate and flips the vertical one:
\[
  \cos(-t) = \cos t \quad(\text{even}),\qquad
  \sin(-t) = -\sin t \quad(\text{odd}).
\]
Cosine is symmetric, sine is antisymmetric. One picture, two identities.

\textbf{Supplementary angle ($\pi - t$).} The point at angle $\pi - t$ is the mirror image of the point at angle $t$ across the \emph{$y$-axis}. Mirroring across the $y$-axis flips $x$ and keeps $y$:
\[
  \cos(\pi - t) = -\cos t,\qquad \sin(\pi - t) = \sin t.
\]

\textbf{Co-function / quarter turn ($\tfrac{\pi}{2} - t$).} This is the deep one, and it is the source of the prefix ``co-'' in cosine. The point at angle $\tfrac{\pi}{2} - t$ is the reflection of the point at angle $t$ across the diagonal line $y = x$. Reflecting across $y=x$ \emph{swaps} the coordinates. Therefore
\[
  \sin\!\left(\tfrac{\pi}{2} - t\right) = \cos t,\qquad
  \cos\!\left(\tfrac{\pi}{2} - t\right) = \sin t.
\]
Cosine of an angle is the sine of its complement -- ``cosine'' literally means ``\textbf{co}mplement's \textbf{sine}.'' That is not a coincidence to memorize; it is the $y=x$ reflection, seen once and never forgotten.

\begin{figure}[H]
\centering
\begin{tikzpicture}[scale=2.5]
  \draw[->] (-1.3,0) -- (1.4,0) node[right]{$x$};
  \draw[->] (0,-1.3) -- (0,1.4) node[above]{$y$};
  \draw[gray!50,thin] (-1.2,-1.2) -- (1.2,1.2) node[right]{\footnotesize $y=x$};
  \draw[curveblue,thick] (0,0) circle (1);
  \def\a{28}
  \filldraw[tangentred] (\a:1) circle (0.6pt) node[below right]{\small $t$};
  \filldraw[structurepurple] ({90-\a}:1) circle (0.6pt) node[above left]{\small $\tfrac{\pi}{2}-t$};
  \draw[thick,tangentred] (0,0) -- (\a:1);
  \draw[thick,structurepurple] (0,0) -- ({90-\a}:1);
  \draw[dashed,gray] (\a:1) -- ({90-\a}:1);
\end{tikzpicture}
\caption{The co-function identity is a mirror. Reflecting the point for angle $t$ across the line $y=x$ gives the point for angle $\tfrac{\pi}{2}-t$, and that reflection swaps the coordinates: $\sin$ and $\cos$ trade places.}
\label{ch04:fig-cofunction}
\end{figure}

\begin{whybox}
Why can we get away with \emph{reading} identities off pictures instead of grinding algebra? Because $\sin$ and $\cos$ were \emph{defined} as coordinates of a point on the circle, and a geometric transformation of the circle (a reflection, a rotation) does a predictable thing to coordinates. Reflect across the $x$-axis: $y \mapsto -y$, so $\sin \mapsto -\sin$. Rotate by a quarter turn, reflect across a diagonal -- each motion is a simple coordinate operation. So every symmetry of the circle is automatically an identity, and the ``identity'' is just the bookkeeping of where the coordinates went. The algebra works because the geometry works; the picture is the proof.
\end{whybox}

\subsection{The Addition Formulas}

The one family of identities you cannot quite \emph{see} in a single glance, but which everything downstream depends on, is the addition formulas:
\[
  \cos(a+b) = \cos a\cos b - \sin a \sin b,
  \qquad
  \sin(a+b) = \sin a\cos b + \cos a\sin b.
\]
These say: if you know where angles $a$ and $b$ land, you can find where angle $a+b$ lands without walking the whole way. They are the engine behind the derivative of sine, so they earn a real derivation. The cleanest one is a rotation argument.

\begin{focusbox}
\textbf{Deriving the addition formula by rotation.} Rotating the plane counterclockwise by an angle $a$ sends the point $(1,0)$ to $(\cos a, \sin a)$ and sends $(0,1)$ to $(-\sin a, \cos a)$ (a quarter turn ahead of the first). A rotation is linear, so it is completely determined by what it does to these two basis points: the point $(x,y)$ goes to
\[
  (x\cos a - y\sin a,\ \ x\sin a + y\cos a).
\]
Now take the point at angle $b$, namely $(x,y) = (\cos b, \sin b)$, and rotate it by $a$. Geometrically that lands it at angle $a+b$, i.e.\ at $(\cos(a+b), \sin(a+b))$. Algebraically the formula above sends it to
\[
  (\cos b\cos a - \sin b\sin a,\ \ \cos b \sin a + \sin b\cos a).
\]
Matching the two coordinates gives exactly the cosine and sine addition formulas. The single idea ``rotating by $a$ then by $b$ is the same as rotating by $a+b$'' contains both identities at once.
\end{focusbox}

Two consequences we will lean on. Setting $b = a$ gives the \textbf{double-angle} formulas:
\[
  \sin 2a = 2\sin a\cos a,\qquad
  \cos 2a = \cos^2 a - \sin^2 a = 1 - 2\sin^2 a = 2\cos^2 a - 1.
\]
The last two forms (using $\cos^2 + \sin^2 = 1$) rearrange into the \textbf{half-angle} formulas, which are workhorses in integration:
\[
  \sin^2 a = \frac{1 - \cos 2a}{2},\qquad
  \cos^2 a = \frac{1 + \cos 2a}{2}.
\]
Keep these last two in your pocket; when you later integrate $\sin^2$ or $\cos^2$, these are how you do it.

\begin{historybox}
The addition formulas are ancient. In effect they appear in Ptolemy's \emph{Almagest} (c.\ 150 CE) as ``Ptolemy's theorem'' about chords of a circle -- Greek astronomers worked not with our sine but with the \emph{chord} of an angle, the straight-line distance across the arc, which equals $2\sin(\theta/2)$ on a unit circle. The word ``sine'' itself is a chain of mistranslations: the Sanskrit \emph{jya} (chord/bowstring) became the Arabic \emph{jiba}, written without vowels as \emph{jb}; later Latin translators misread \emph{jb} as \emph{jaib}, meaning ``bay'' or ``fold,'' and rendered it \emph{sinus} -- Latin for a bay or the fold of a garment. So every time you write $\sin$, you are writing the Latin for ``bay,'' by way of a twelfth-century reading error. The mathematics is older and clearer than its name.
\end{historybox}

\section{The Graphs: Unrolling the Circle}
\label{ch04:graphs}

We have been watching a point go around a circle. To graph $\sin t$ and $\cos t$ as functions of $t$, we \emph{unroll} that circular motion: put $t$ on the horizontal axis (how far around we have walked) and the coordinate on the vertical axis. The circle's repetition becomes a wave that rolls on forever.

\begin{figure}[H]
\centering
\begin{tikzpicture}
\begin{axis}[
  width=13cm, height=6cm,
  xmin=-0.4, xmax=12.8, ymin=-1.5, ymax=1.5,
  axis lines=middle,
  xtick={0,1.5708,3.1416,4.7124,6.2832,7.854,9.4248,10.996,12.566},
  xticklabels={$0$,$\frac{\pi}{2}$,$\pi$,$\frac{3\pi}{2}$,$2\pi$,$\frac{5\pi}{2}$,$3\pi$,$\frac{7\pi}{2}$,$4\pi$},
  ytick={-1,1},
  xlabel={$t$}, ylabel={},
  samples=200, domain=0:12.566,
  legend pos=south east, legend style={font=\small},
]
  \addplot[curveblue,very thick] {sin(deg(x))};
  \addlegendentry{$\sin t$}
  \addplot[tangentred,very thick,dashed] {cos(deg(x))};
  \addlegendentry{$\cos t$}
  \draw[gray!50,dashed] (axis cs:0,1) -- (axis cs:12.566,1);
  \draw[gray!50,dashed] (axis cs:0,-1) -- (axis cs:12.566,-1);
\end{axis}
\end{tikzpicture}
\caption{Sine and cosine, the unrolled shadows of circular motion. Each is a wave of amplitude $1$ and period $2\pi$. Cosine is sine shifted left by $\pi/2$ -- the co-function identity, made visual.}
\label{ch04:fig-sincos-graph}
\end{figure}

Read the structure straight off the graph. The wave oscillates between $-1$ and $1$ (the \textbf{amplitude} is $1$). It completes one full cycle every $2\pi$ (the \textbf{period}). Sine starts at $0$ climbing; cosine starts at $1$ already at its peak. And the two curves are \emph{the same shape shifted}: $\cos t = \sin(t + \tfrac\pi2)$. Shifting sine to the left by a quarter period lands it exactly on cosine. (This is the co-function identity in motion: using $\sin$ odd, $\sin(t+\tfrac\pi2) = \sin(\tfrac\pi2-(-t)) = \cos(-t) = \cos t$.) More simply still: the cosine wave is the sine wave with a head start of a quarter turn, which is precisely what a quarter turn around the circle does.

\begin{connectionbox}
A general sinusoid is written $y = A\sin(\omega t + \varphi) + k$. Each parameter is a geometric knob: $A$ is the amplitude (vertical stretch), $\omega$ is the angular frequency (horizontal compression -- the period is $2\pi/\omega$), $\varphi$ is the phase (horizontal shift, the head start), and $k$ is the vertical offset (the midline). Because sine and cosine differ only by a phase, \emph{any} sinusoid can be written with either one. This four-knob family is exactly what describes simple harmonic motion in physics -- a mass on a spring, $x(t) = A\cos(\omega t + \varphi)$ -- and it is the building block of Fourier analysis, where astonishingly \emph{every} periodic signal, however jagged, is a sum of these smooth waves.
\end{connectionbox}

The tangent graph looks entirely different, and the circle tells us why. As $t$ climbs toward $\pi/2$, the point nears the top of the circle: $\sin t \to 1$ while $\cos t \to 0^+$, so $\tan t = \sin t/\cos t \to +\infty$. Just past $\pi/2$, cosine is negative and tangent leaps to $-\infty$. The result is a curve that races up to a vertical asymptote, vanishes off the top, and reappears from the bottom -- repeating every $\pi$, not $2\pi$.

\begin{figure}[H]
\centering
\begin{tikzpicture}
\begin{axis}[
  width=12cm, height=6cm,
  xmin=-5, xmax=5, ymin=-5, ymax=5,
  axis lines=middle,
  xtick={-4.7124,-1.5708,1.5708,4.7124},
  xticklabels={$-\frac{3\pi}{2}$,$-\frac{\pi}{2}$,$\frac{\pi}{2}$,$\frac{3\pi}{2}$},
  ytick={-4,-2,2,4},
  xlabel={$t$},
  restrict y to domain=-6:6,
  samples=200,
]
  \addplot[curveblue,very thick,domain=-1.45:1.45] {tan(deg(x))};
  \addplot[curveblue,very thick,domain=1.69:4.59] {tan(deg(x))};
  \addplot[curveblue,very thick,domain=-4.59:-1.69] {tan(deg(x))};
  \draw[gray!55,dashed] (axis cs:1.5708,-6) -- (axis cs:1.5708,6);
  \draw[gray!55,dashed] (axis cs:-1.5708,-6) -- (axis cs:-1.5708,6);
  \draw[gray!55,dashed] (axis cs:4.7124,-6) -- (axis cs:4.7124,6);
  \draw[gray!55,dashed] (axis cs:-4.7124,-6) -- (axis cs:-4.7124,6);
\end{axis}
\end{tikzpicture}
\caption{The tangent function $\tan t = \sin t/\cos t$. Vertical asymptotes occur where $\cos t = 0$. The function runs from $-\infty$ to $+\infty$ on each interval of width $\pi$, then repeats.}
\label{ch04:fig-tan-graph}
\end{figure}

\section{The Small-Angle Behavior: $\sin t \approx t$}
\label{ch04:smallangle}

This section hands the rest of the book its single most important fact about trig. Look again at the unit circle for a \emph{small} angle $t$, just above $(1,0)$.

\begin{figure}[H]
\centering
\begin{tikzpicture}[scale=4.2]
  \draw[->] (-0.1,0) -- (1.25,0) node[right]{$x$};
  \draw[->] (0,-0.05) -- (0,0.62) node[above]{$y$};
  \draw[curveblue,thick] (0,0) arc (0:38:1) ;
  \def\a{22}
  \draw[thick] (0,0) -- (\a:1);
  \filldraw (\a:1) circle (0.3pt);
  % the arc t (thick red)
  \draw[tangentred,very thick] (1,0) arc (0:\a:1);
  % sin t : vertical from point to x axis
  \draw[structurepurple,very thick] (\a:1) -- ({cos(\a)},0);
  % tan t: vertical line at x=1 up to extended radius
  \draw[targetgreen,very thick] (1,0) -- (1,{tan(\a)});
  \draw[thick] (\a:1) -- ({1},{tan(\a)});
  \node[structurepurple] at ({cos(\a)+0.05},{sin(\a)/2}) {\small $\sin t$};
  \node[tangentred] at ({1.07},{sin(\a)/2 - 0.02}) {\small $t$};
  \node[targetgreen] at ({1.13},{tan(\a)/1.5}) {\small $\tan t$};
  \node at ({cos(\a)},-0.04) {\small $\cos t$};
\end{tikzpicture}
\caption{For a small angle $t$, the vertical segment $\sin t$, the arc $t$, and the tangent segment $\tan t$ are nearly equal, with $\sin t < t < \tan t$. As $t\to 0$ the three become indistinguishable, which forces $\sin t/t \to 1$.}
\label{ch04:fig-smallangle}
\end{figure}

Three lengths are racing to zero together as $t\to 0$: the vertical drop $\sin t$, the arc $t$ along the rim, and the tangent segment $\tan t$. Geometrically they nest:
\[
  \sin t \;<\; t \;<\; \tan t \qquad (0 < t < \tfrac{\pi}{2}).
\]
The first inequality says the straight vertical chord is shorter than the curved arc above it; the second compares two triangles against the sector. As the angle shrinks, the arc straightens out and these three lengths become indistinguishable. The most useful way to say this precisely is in terms of the ratio $\sin t / t$.

\begin{focusbox}
\textbf{The fundamental trig limit.}
\[
  \boxed{\;\lim_{t\to 0}\frac{\sin t}{t} = 1.\;}
\]
In words: for small $t$ (\emph{in radians}), $\sin t$ and $t$ are almost equal -- their ratio rushes to $1$. Equivalently, $\sin t \approx t$ when $t$ is small -- Figure~\ref{ch04:fig-sint-approx-t} shows the curve and the line clinging together near the origin. This is the seed from which the entire calculus of trig functions grows; the derivative of sine, as you will see in the next chapter, is exactly this limit in disguise.
\end{focusbox}

We can already see \emph{why} it must be true, and even sketch the rigorous proof, using only the sector-area formula $A = \tfrac12 r^2\theta$ from the start of the chapter.

\begin{focusbox}
\textbf{Why $\sin t / t \to 1$: the squeeze.} Work on the unit circle ($r=1$) with $0 < t < \pi/2$. Compare three areas in Figure~\ref{ch04:fig-smallangle}:
\begin{itemize}
  \item the inner triangle with vertices $(0,0),(1,0),(\cos t,\sin t)$ has area $\tfrac12\sin t$;
  \item the circular sector of angle $t$ has area $\tfrac12 t$ (using $A=\tfrac12 r^2\theta$ with $r=1$);
  \item the outer triangle with the vertical tangent, vertices $(0,0),(1,0),(1,\tan t)$, has area $\tfrac12\tan t$.
\end{itemize}
The inner triangle sits inside the sector, which sits inside the outer triangle, so
\[
  \tfrac12 \sin t \;<\; \tfrac12 t \;<\; \tfrac12 \tan t.
\]
Multiply by $2$ and divide by $\sin t$ (positive here):
\[
  1 \;<\; \frac{t}{\sin t} \;<\; \frac{1}{\cos t}.
\]
Flip everything (taking reciprocals reverses the inequalities):
\[
  \cos t \;<\; \frac{\sin t}{t} \;<\; 1.
\]
As $t\to 0$, $\cos t \to 1$. The ratio $\sin t/t$ is trapped -- \emph{squeezed} -- between something heading to $1$ and the constant $1$. It has no choice: $\lim_{t\to 0}\sin t/t = 1$. (The case $t<0$ follows because $\sin t/t$ is even.)
\end{focusbox}

\begin{trapbox}
This limit equals $1$ \textbf{only in radians}. If $t$ is measured in degrees, then walking ``$t$ degrees'' covers an arc of only $\tfrac{\pi}{180}t$, so $\sin(t^\circ)/t \to \pi/180 \approx 0.01745$, a needless ugly constant. Every derivative formula for trig functions would be polluted by that factor forever. \emph{This is the real reason calculus uses radians.} Radians are the unit in which arc length equals angle, which is the unit in which $\sin t \approx t$, which is the unit in which the derivatives come out clean.
\end{trapbox}

\begin{figure}[H]
\centering
\begin{tikzpicture}
\begin{axis}[
  width=11cm, height=6cm,
  xmin=-1.6, xmax=1.6, ymin=-1.6, ymax=1.6,
  axis lines=middle,
  xtick={-1.5708,-0.7854,0.7854,1.5708},
  xticklabels={$-\frac{\pi}{2}$,$-\frac{\pi}{4}$,$\frac{\pi}{4}$,$\frac{\pi}{2}$},
  ytick={-1,1},
  xlabel={$t$}, ylabel={},
  samples=120, domain=-1.5708:1.5708,
  legend pos=south east, legend style={font=\small},
]
  \addplot[tangentred,very thick,dashed] {x};
  \addlegendentry{$y=t$}
  \addplot[curveblue,very thick] {sin(deg(x))};
  \addlegendentry{$\sin t$}
\end{axis}
\end{tikzpicture}
\caption{Near $t=0$ the curve $\sin t$ hugs the line $y=t$: they share the same value and the same slope $1$ at the origin, and only peel apart slowly. This is the visual content of $\sin t/t\to 1$, equivalently $\sin t\approx t$ for small $t$.}
\label{ch04:fig-sint-approx-t}
\end{figure}

\begin{aestheticbox}
Pause on how much rides on $\sin t/t\to 1$. It is the statement that, up close, a sine wave and a straight line of slope $1$ are the same thing -- that the circle, examined under a microscope at $(1,0)$, looks flat and vertical, rising one unit of height per one unit of arc. From this single ``the curve is locally straight'' fact will come $\dv{}{t}\sin t = \cos t$, the wave equation, the description of light and sound, and the Taylor series $\sin t = t - \tfrac{t^3}{6} + \cdots$ whose very first term is this approximation $\sin t\approx t$. A great deal of physics is the sentence ``for small angles, $\sin t \approx t$'' taken seriously.
\end{aestheticbox}

\begin{applicationbox}
\textbf{The pendulum.} A pendulum of length $L$ obeys, exactly, $\ddot\theta = -\tfrac{g}{L}\sin\theta$ -- a hard, nonlinear equation with no elementary solution. But for small swings, $\sin\theta \approx \theta$, and the equation \emph{linearizes} to $\ddot\theta = -\tfrac{g}{L}\theta$, whose solution is a pure sinusoid with period $T = 2\pi\sqrt{L/g}$. That is why a pendulum clock keeps near-constant time regardless of how far it swings (as long as the swing is small): Galileo's observation, explained in one approximation. The whole of small-oscillation physics -- springs, molecules, circuits, bridges -- rests on replacing a curve by its tangent line, and $\sin\theta\approx\theta$ is the first and most famous instance.
\end{applicationbox}

A companion limit, derived the same way and needed alongside the first, is
\[
  \lim_{t\to 0}\frac{1 - \cos t}{t} = 0.
\]
Intuitively: near $t=0$ the height $\sin t$ falls off the circle like $t$ (first order), but the horizontal gap $1-\cos t$ falls off like $t^2/2$ (second order, much faster), so dividing it by $t$ still leaves something heading to zero. We prove it with a clean algebraic trick -- multiply by the conjugate:
\[
  \frac{1-\cos t}{t}
  = \frac{(1-\cos t)(1+\cos t)}{t(1+\cos t)}
  = \frac{1-\cos^2 t}{t(1+\cos t)}
  = \frac{\sin^2 t}{t(1+\cos t)}
  = \underbrace{\frac{\sin t}{t}}_{\to 1}\cdot \underbrace{\frac{\sin t}{1+\cos t}}_{\to\, 0/2\,=\,0}.
\]
The first factor tends to $1$ by our main limit; the second tends to $0/2 = 0$. Product: $0$. Both limits will reappear, fused, the moment we differentiate sine.

\section{Inverse Trigonometric Functions}
\label{ch04:inverse}

We turn now from ``given the angle, find the coordinate'' to its reverse: ``given the coordinate, find the angle.'' If you know that a point on the unit circle has height $\sin t = 0.5$, which $t$ produced it? This is the job of the inverse trig functions -- but to define them properly we must first confront a real obstacle.

\subsection{The Obstacle: Periodic Functions Are Not One-to-One}

A function has an inverse only if it is \textbf{one-to-one}: each output comes from exactly one input. Graphically, this is the horizontal line test -- no horizontal line may cross the graph more than once, or else that output traces back to several inputs and ``the inverse'' is ambiguous.

Sine fails this test spectacularly. The horizontal line $y = 0.5$ slices through the sine graph infinitely often: $\sin t = 0.5$ at $t = \pi/6$, at $t = 5\pi/6$, at $t = \pi/6 + 2\pi$, and at infinitely many other places. ``The angle whose sine is $0.5$'' has no single answer. A periodic function repeats, so by its nature it can never be globally invertible.

\begin{intuitionbox}
The fix is blunt and a little ruthless: we \emph{restrict} sine to a stretch of its domain on which it never repeats a value -- a stretch where it is strictly increasing (or strictly decreasing) and sweeps through its full range exactly once. On that restricted piece the inverse exists and is unique. We are not inverting all of sine -- that is impossible. We are inverting one carefully chosen, representative slice of it, and then living forever with the consequences of that choice.
\end{intuitionbox}

\subsection{Arcsine: Choosing the Principal Branch}

For sine, the agreed-upon slice is the interval $[-\tfrac{\pi}{2}, \tfrac{\pi}{2}]$. Why this one? On it, sine climbs monotonically from $\sin(-\tfrac\pi2) = -1$ up to $\sin(\tfrac\pi2) = 1$, hitting every value in $[-1,1]$ exactly once. It is the slice straddling the origin, it is symmetric, and it captures sine at its most natural -- rising through zero with slope $1$. This restricted, invertible sine has an inverse called the \textbf{arcsine}.

\begin{definition}[Arcsine]
\label{ch04:def-arcsin}
For $-1 \le x \le 1$, $\arcsin x$ (also written $\sin^{-1}x$) is the unique angle $\theta \in [-\tfrac{\pi}{2}, \tfrac{\pi}{2}]$ with $\sin\theta = x$. Its domain is $[-1,1]$ (the possible sine values) and its range is $[-\tfrac{\pi}{2}, \tfrac{\pi}{2}]$ (the chosen slice of angles). The name \emph{arc}sine reflects exactly what it returns: the \emph{arc} (angle) belonging to a given coordinate.
\end{definition}

\begin{trapbox}
The superscript notation $\sin^{-1}x$ is a genuine menace. It does \textbf{not} mean $1/\sin x$ -- that is $\csc x$ or $(\sin x)^{-1}$. Here the ``$-1$'' marks \emph{functional inverse}, not a reciprocal power. The clash with $\sin^2 x = (\sin x)^2$ is an unfortunate historical inconsistency in the same notation. When in doubt, write $\arcsin x$; it cannot be misread.
\end{trapbox}

The graph of an inverse function is the \emph{reflection} of the original across the line $y=x$ -- because inverting a function swaps the roles of input and output, which swaps the axes. So $\arcsin$ is the restricted sine curve flipped over the diagonal.

\begin{figure}[H]
\centering
\begin{tikzpicture}
\begin{axis}[
  width=8.5cm, height=8.5cm,
  axis equal image,
  xmin=-1.9, xmax=1.9, ymin=-1.9, ymax=1.9,
  axis lines=middle,
  xtick={-1.5708,-1,1,1.5708}, xticklabels={$-\frac{\pi}{2}$,$-1$,$1$,$\frac{\pi}{2}$},
  ytick={-1.5708,-1,1,1.5708}, yticklabels={$-\frac{\pi}{2}$,$-1$,$1$,$\frac{\pi}{2}$},
  samples=120,
  legend pos=north west, legend style={font=\footnotesize},
]
  \addplot[gray!55,dashed,domain=-1.9:1.9] {x};
  \addlegendentry{$y=x$}
  \addplot[curveblue,very thick,domain=-1.5708:1.5708] {sin(deg(x))};
  \addlegendentry{$\sin t$ (restricted)}
  \addplot[tangentred,very thick,domain=-1:1] {deg(asin(x))*pi/180};
  \addlegendentry{$\arcsin x$}
\end{axis}
\end{tikzpicture}
\caption{Restricted sine (on $[-\tfrac\pi2,\tfrac\pi2]$) and its mirror image across $y=x$, the arcsine. Reflecting a graph across $y=x$ is what inverting a function does: it swaps input and output.}
\label{ch04:fig-arcsin}
\end{figure}

\subsection{Arccosine and Arctangent}

The same surgery defines the other inverses; only the chosen slice differs, each picked to be a monotonic stretch covering the full range once.

\textbf{Arccosine.} Cosine \emph{decreases} from $1$ to $-1$ on $[0,\pi]$, sweeping its whole range once, so that is the slice. Define $\arccos x$ for $x\in[-1,1]$ as the unique $\theta\in[0,\pi]$ with $\cos\theta = x$. Note its range starts at $0$, not $-\tfrac\pi2$: a different function needs a different home interval. A pretty relation ties the two together. Since $\cos\theta = \sin(\tfrac\pi2 - \theta)$ (the co-function identity), one can show
\[
  \arcsin x + \arccos x = \frac{\pi}{2}\qquad\text{for all }x\in[-1,1].
\]
The arcsine and arccosine of the same number are complementary angles -- again the ``co-'' relationship, now read backwards through the inverses.

\textbf{Arctangent.} Tangent runs monotonically from $-\infty$ to $+\infty$ on the open interval $(-\tfrac\pi2, \tfrac\pi2)$, so that is its slice. Define $\arctan x$ for \emph{every} real $x$ as the unique $\theta\in(-\tfrac\pi2,\tfrac\pi2)$ with $\tan\theta = x$. Because tangent's range is all of $\R$, arctangent accepts every real input -- a genuinely useful feature. Its graph has two horizontal asymptotes, at $y = \pm\tfrac\pi2$: feed in an enormous number and the angle approaches, but never reaches, a quarter turn.

\begin{figure}[H]
\centering
\begin{tikzpicture}
\begin{axis}[
  width=12cm, height=5.4cm,
  xmin=-8, xmax=8, ymin=-2, ymax=2,
  axis lines=middle,
  ytick={-1.5708,1.5708}, yticklabels={$-\frac{\pi}{2}$,$\frac{\pi}{2}$},
  xlabel={$x$},
  samples=160, domain=-8:8,
]
  \addplot[curveblue,very thick] {deg(atan(x))*pi/180};
  \draw[gray!55,dashed] (axis cs:-8,1.5708) -- (axis cs:8,1.5708);
  \draw[gray!55,dashed] (axis cs:-8,-1.5708) -- (axis cs:8,-1.5708);
\end{axis}
\end{tikzpicture}
\caption{Arctangent: defined for all real $x$, with horizontal asymptotes at $\pm\pi/2$. It is the smooth, bounded ``angle-of'' function -- which is why it appears everywhere a saturating or angle-recovering response is needed.}
\label{ch04:fig-arctan}
\end{figure}

\begin{center}
\begin{tabular}{lccc}
\toprule
Inverse & Domain (inputs) & Range (outputs) & Chosen slice of the original \\
\midrule
$\arcsin x$ & $[-1,1]$ & $[-\tfrac\pi2,\tfrac\pi2]$ & sine increasing through $0$ \\
$\arccos x$ & $[-1,1]$ & $[0,\pi]$ & cosine decreasing from $1$ \\
$\arctan x$ & $(-\infty,\infty)$ & $(-\tfrac\pi2,\tfrac\pi2)$ & tangent increasing through $0$ \\
\bottomrule
\end{tabular}
\end{center}

\begin{happeningbox}
What an inverse trig function ``really'' returns is a \emph{location on the circle}, reported as the arc length from $(1,0)$ -- but \emph{only the location within the chosen branch}. When your calculator says $\arcsin(0.5) = \pi/6$, it is not lying, but it is silently making a choice: it reports the one answer in $[-\tfrac\pi2,\tfrac\pi2]$ and quietly discards the equally-valid $5\pi/6$ and all the $+2\pi k$ cousins. This is why solving $\sin t = 0.5$ in full requires you to add back what the inverse threw away: $t = \pi/6 + 2\pi k$ \emph{and} $t = 5\pi/6 + 2\pi k$. The inverse gives you one root; the symmetries of the circle give you the rest.
\end{happeningbox}

\begin{whybox}
Why must we restrict the domain at all, rather than just ``allowing'' a multivalued inverse? Because in calculus we need a genuine \emph{function} -- one output per input -- so that we can graph it, differentiate it, and feed it into the chain rule. A function that returns infinitely many answers cannot have a single well-defined derivative at a point. By amputating sine down to one monotone branch we buy a clean, differentiable, single-valued inverse. The cost is that the inverse only ever reports angles in one limited window; we accept that cost because differentiability is worth more than completeness, and because the discarded answers are always recoverable from the periodicity and symmetry of the original.
\end{whybox}

\subsection{Computing With Inverses: the Triangle Trick}

A frequent and initially baffling task is to simplify a composition like $\cos(\arcsin x)$. The expression mixes an angle-finder with a coordinate-finder, but a right-triangle picture untangles it instantly.

\begin{example}[Simplify $\cos(\arcsin x)$]
\label{ch04:ex-triangle}
Let $\theta = \arcsin x$, so $\sin\theta = x$ and, by definition of the arcsine branch, $\theta\in[-\tfrac\pi2,\tfrac\pi2]$. We want $\cos\theta$. Build a right triangle in which $\sin\theta = x$: opposite side $x$, hypotenuse $1$ (so that opposite/hypotenuse $= x$). By Pythagoras the adjacent side is $\sqrt{1 - x^2}$. Then
\[
  \cos(\arcsin x) = \cos\theta = \frac{\text{adjacent}}{\text{hypotenuse}} = \sqrt{1 - x^2}.
\]
We take the \emph{positive} root because $\theta$ lies in $[-\tfrac\pi2,\tfrac\pi2]$, where cosine is never negative. The branch we chose for arcsine is exactly what resolves the sign ambiguity -- another reason the restriction matters.
\end{example}

This little identity, $\cos(\arcsin x) = \sqrt{1-x^2}$, is not a curiosity: it is precisely what makes the derivative of $\arcsin$ come out to the clean $1/\sqrt{1-x^2}$ you will derive next chapter, and it is the substitution behind a whole family of integrals.

\begin{labbox}
Try the same trick on $\tan(\arccos x)$. Set $\theta=\arccos x$, so $\cos\theta = x$ with adjacent $=x$, hypotenuse $=1$, hence opposite $=\sqrt{1-x^2}$. Then $\tan\theta = \text{opposite}/\text{adjacent} = \sqrt{1-x^2}/x$. Sketch the triangle yourself and confirm each side; the picture does all the work.
\end{labbox}

\section{Why Radians Make Calculus Clean: The Payoff}
\label{ch04:why-radians}

One thread has run through the whole chapter, and it is the most important practical takeaway. Why radians?

Trace the logic. We \emph{defined} radians so that arc length equals angle ($s = r\theta$, and $\theta$ itself on the unit circle). That choice made the arc above $(1,0)$ have length exactly $t$. That, in turn, forced the squeeze $\cos t < \tfrac{\sin t}{t} < 1$, hence $\sin t/t \to 1$, hence $\sin t \approx t$ for small $t$. And \emph{that} approximation is, almost verbatim, the statement that the derivative of $\sin t$ at $t=0$ equals $1$:
\[
  \dv{}{t}\sin t \Big|_{t=0} = \lim_{t\to 0}\frac{\sin t - \sin 0}{t - 0} = \lim_{t\to 0}\frac{\sin t}{t} = 1.
\]
With that, the next chapter's central results $\dv{}{t}\sin t = \cos t$ and $\dv{}{t}\cos t = -\sin t$ fall out clean, with no stray constants. Had we used degrees, every one of these would carry a factor of $\pi/180 \approx 0.01745$, multiplying and compounding through every higher derivative, every integral, every Taylor series. Radians are not a convention chosen for tidiness; they are the \emph{unique} angle unit in which the calculus of trigonometry has no extraneous constant at all.

\begin{aestheticbox}
This is a recurring pattern in mathematics worth naming. The right \emph{units} or the right \emph{definitions} do not merely simplify the bookkeeping -- they reveal that two things you thought were different are secretly the same. Radians reveal that angle and arc length are one quantity. The number $e$ (you have met it, or will) is the base that reveals that an exponential and its own rate of growth are one. In both cases, choosing the ``natural'' unit makes a derivative equal to $1$ exactly where it should, and the whole subject snaps into focus. When the formulas suddenly lose their clutter, it is usually a sign you have found the coordinates the problem wanted all along.
\end{aestheticbox}

\section{Exercises}
\label{ch04:exercises}

\begin{exercise}[Conversions and arc length]
(a) Convert to radians: $135^\circ$, $210^\circ$, $300^\circ$.
(b) Convert to degrees: $\tfrac{5\pi}{6}$, $\tfrac{7\pi}{4}$, $3$ radians (approximate).
(c) A circular arc has radius $4$ cm and subtends an angle of $\tfrac{2\pi}{3}$. Find its length and the area of the corresponding sector.
\end{exercise}

\begin{exercise}[Reading the circle]
Without a calculator, give exact values of $\sin$, $\cos$, and $\tan$ at $t = \tfrac{2\pi}{3}$, $t = \tfrac{5\pi}{4}$, and $t = \tfrac{11\pi}{6}$. For each, state which quadrant the point lands in and use that to predict the signs \emph{before} computing.
\end{exercise}

\begin{exercise}[Identities from symmetry]
Using only the unit-circle picture (reflections and rotations), explain why each holds:
(a) $\sin(\pi + t) = -\sin t$;\quad
(b) $\cos(\pi + t) = -\cos t$;\quad
(c) $\tan(\pi + t) = \tan t$;\quad
(d) $\sin(\tfrac{\pi}{2} + t) = \cos t$. Draw the relevant reflection or rotation for at least two of them.
\end{exercise}

\begin{exercise}[Double and half angle]
(a) Starting from the addition formula for $\cos(a+b)$, derive $\cos 2a = 1 - 2\sin^2 a$.
(b) Use the half-angle formula to find the exact value of $\cos^2(\pi/8)$, and hence $\cos(\pi/8)$.
\end{exercise}

\begin{exercise}[The fundamental limit at work]
Evaluate each limit, reducing it to $\lim_{t\to0}\sin t/t = 1$ (and its cousin $\tfrac{1-\cos t}{t}\to 0$):
(a) $\displaystyle\lim_{t\to 0}\frac{\sin 5t}{t}$;\quad
(b) $\displaystyle\lim_{t\to 0}\frac{\sin 3t}{\sin 7t}$;\quad
(c) $\displaystyle\lim_{t\to 0}\frac{\tan t}{t}$;\quad
(d) $\displaystyle\lim_{t\to 0}\frac{1 - \cos t}{t^2}$.
(Hint for (a): write $\tfrac{\sin 5t}{t} = 5\cdot\tfrac{\sin 5t}{5t}$.)
\end{exercise}

\begin{exercise}[Inverse compositions]
Simplify to an algebraic expression in $x$ (state any restrictions on $x$):
(a) $\sin(\arccos x)$;\quad
(b) $\tan(\arcsin x)$;\quad
(c) $\cos(\arctan x)$;\quad
(d) $\sin(2\arctan x)$.
\end{exercise}

\begin{exercise}[Why the branch matters]
(a) Is it true that $\arcsin(\sin t) = t$ for \emph{every} real $t$? Test $t = \tfrac{3\pi}{4}$ and explain the result using the range of $\arcsin$.
(b) For which $t$ \emph{is} $\arcsin(\sin t) = t$ exactly true?
(c) Sketch the function $g(t) = \arcsin(\sin t)$ over $[-2\pi, 2\pi]$ and describe its shape.
\end{exercise}

\begin{exercise}[Solving trig equations completely]
Find \emph{all} real solutions, using an inverse for one solution and the circle's symmetry for the rest:
(a) $2\sin t - 1 = 0$;\quad
(b) $\cos t = -\tfrac{\sqrt3}{2}$;\quad
(c) $\tan t = 1$. Express each answer as a family $t = (\ldots) + 2\pi k$ or $+\pi k$, $k\in\Z$.
\end{exercise}

\begin{exercise}[Prove the squeeze, carefully]
Reproduce the proof that $\lim_{t\to 0^+}\sin t/t = 1$ by comparing the areas of the inner triangle, the sector, and the outer triangle on the unit circle. Then explain why the two-sided limit follows, by arguing that $f(t) = \sin t/t$ is an even function.
\end{exercise}

\begin{exercise}[A modeling problem]
The number of daylight hours at a mid-northern latitude is modeled by
\[
  D(t) = 12 + 3\sin\!\left(\frac{2\pi}{365}(t - 80)\right),
\]
where $t$ is the day of the year. (a) What are the amplitude, period, and midline, and what do they mean physically? (b) On which day is daylight longest? (c) For roughly how many days per year does $D(t) > 13.5$? (Use an inverse sine and the symmetry of the curve.)
\end{exercise}

\begin{exercise}[Conceptual: the missing constant]
Suppose a careless textbook \emph{defines} sine using degrees, so that $S(x) = \sin(x^\circ)$ where $x$ is a pure number of degrees. Using $\sin(x^\circ) = \sin\!\big(\tfrac{\pi}{180}x \text{ rad}\big)$ and the radian limit, compute $\lim_{x\to0} S(x)/x$. Explain in one or two sentences why this answer makes degrees a poor choice for calculus, connecting it to the eventual formula for the derivative of $S$.
\end{exercise}
